\clearpage
\section{Technologies Used}

\newcommand{\techitem}[3]{%
  \subsection{#1}
  \begin{minipage}[c]{0.70\textwidth}#2\end{minipage}\hfill
  \begin{minipage}[c]{0.24\textwidth}\centering\includegraphics[width=0.85\linewidth]{logos/#3}\end{minipage}
  \vspace{1.2em}}

\techitem{Docker}{Docker is a containerization platform that allows packaging an application
and its dependencies into a lightweight, portable container. A multi-stage Dockerfile builds
the React frontend and the Python API into a single image, which standardizes the development
environment and simplifies the deployment of the API and web application.}{docker.png}

\techitem{React}{React is an open-source JavaScript library developed by Meta, used to create
interactive user interfaces. It was used (with Vite) to develop the \textbf{Frontend} of our
application, providing a smooth and responsive user experience for composition input,
3D crystal visualization (3Dmol.js) and prediction display.}{react.png}

\techitem{FastAPI}{FastAPI is a modern and high-performance web framework for building APIs
with Python. It was chosen as the \textbf{Backend} to serve our Machine Learning models due to
its speed, native asynchronous support, and automatic documentation generation.}{fastapi.png}

\techitem{Jupyter Notebook}{Jupyter Notebook is an open-source web application for creating
and sharing documents containing code, equations, and visualizations. It served as the main
tool for data enhancement, Exploratory Data Analysis (EDA), preprocessing, and the
presentation of the SSL experiments.}{jupyter.png}

\techitem{GitHub}{GitHub is a Git-based source code management platform. It was used for
project hosting, team collaboration, continuous integration (GitHub Actions: linting and
tests) and version control of the source code.}{github.png}

\techitem{pymatgen \& matminer}{pymatgen and matminer are the reference open-source Python
libraries of materials informatics. They were essential for parsing chemical formulas,
handling compositions and crystal structures, downloading the Matbench datasets and defining
the Magpie descriptors needed for the models.}{pymatgen.png}

\subsection{scikit-learn, LightGBM and Optuna}
scikit-learn provides the semi-supervised algorithms (Self-Training, Label Propagation, Label
Spreading), the classical models and the evaluation metrics; LightGBM and XGBoost provide fast
gradient-boosting learners, and Optuna performs the Bayesian (TPE) hyper-parameter search used
to tune the base learners. DEAP implements the genetic algorithm and mlxtend the Apriori
association-rule mining.
